% Artículo del Proyecto 1: formato IEEE Transactions (A1, A2), máximo 4 páginas (A3).
% Sigue la plantilla oficial de IEEE para Transactions y Journals
% (IEEE-Transactions-LaTeX2e-templates-and-instructions.zip, IEEE Template
% Selector, versión del 24/10/2025): bare_jrnl_new_sample4.tex y
% New_IEEEtran_how-to.pdf. IEEEtran.cls en esta carpeta es la de esa plantilla.
%
% Reglas de la plantilla que aplican al texto:
% - El primer párrafo de la introducción empieza con \IEEEPARstart.
% - Figuras con [!t], \includegraphics antes de \caption y \label después
%   de \caption.
% - Tablas con \caption antes de la tabla y \label después. Las divisiones
%   las fija el enunciado (A10, A11): solo horizontales, con booktabs.
% - Ecuaciones numeradas con equation o align (no eqnarray ni $$), citadas con
%   \eqref; sin fórmulas en el título ni en el resumen.
% - Las referencias se citan con \cite, que el paquete cite ordena y agrupa.
% La muestra oficial usa la opción 'lettersize', que su IEEEtran.cls no define
% (la ignora con una advertencia); 'letterpaper' es la opción equivalente que sí existe.
\documentclass[letterpaper,journal]{IEEEtran}

\usepackage[T1]{fontenc}
\usepackage[utf8]{inputenc}
% es-tabla: "Tabla" en lugar de "Cuadro" (A7); es-nodecimaldot: punto decimal.
\usepackage[english,spanish,es-tabla,es-nodecimaldot]{babel}
\usepackage{amsmath,amsfonts,amssymb}
\usepackage{cite}
\usepackage{graphicx}
\usepackage{booktabs} % tablas sin divisiones verticales (A10, A11)
\usepackage{siunitx}
\usepackage[hidelinks]{hyperref}
\usepackage{cleveref}

\graphicspath{{../figuras/}}
\sisetup{output-decimal-marker={.}, per-mode=symbol}
\renewcommand{\IEEEkeywordsname}{Palabras clave}

\begin{document}

\title{Esteganografía de metadatos en audio mediante enmascaramiento con eco}

\author{%
  Marco Zolla Salazar,
  Integrante~A % TODO: nombre (rol A); cada integrante cambia solo su renglón
  y~Integrante~B% TODO: nombre (rol B)
  \thanks{Maestría en Electrónica, Instituto Tecnológico de Costa Rica.}}

% Encabezados de página: revista (aquí, el curso) y autores: título.
\markboth{MP6104 Procesamiento Digital de Señales, Proyecto 1, octubre~2026}%
{Zolla Salazar \MakeLowercase{\textit{et al.}}: Esteganografía de metadatos en audio mediante enmascaramiento con eco}

\maketitle

% Dueño: C. A14: el resumen va en inglés. A15: debe permitir decidir si el artículo es útil.
\begin{otherlanguage}{english}
\begin{abstract}
TODO.
\end{abstract}
\end{otherlanguage}

% Dueño: C. A13: en orden alfabético. A16: representativas del contenido.
\begin{IEEEkeywords}
TODO.
\end{IEEEkeywords}

% Cada archivo de secciones/ tiene un solo dueño (anotado en su primera línea);
% C, como editor, revisa el estilo y la forma (A4-A14).
\section{Introducción}
% Dueño: B. A17: contexto, trabajos relacionados y estado del arte. A18: termina con la estructura del documento.
\IEEEPARstart{L}{a} esteganografía en señales de audio consiste en ocultar información digital dentro de una señal portadora acústica sin alterar de manera perceptible su reproducción, lo que resulta de gran utilidad en el almacenamiento de metadatos de canciones, el control de derechos de autor y la autenticación de grabaciones~\cite{dutta2009,kaur2014}.

Entre las técnicas clásicas reportadas en el estado del arte destacan la sustitución del bit menos significativo, la codificación de fase, el espectro disperso y el enmascaramiento con eco~\cite{dutta2009,kaur2014}. El enmascaramiento con eco, introducido originalmente por Gruhl, Lu y Bender~\cite{gruhl_echo_hiding}, aprovecha la propiedad psicoacústica según la cual una réplica atenuada con un retardo del orden de un milisegundo no se percibe como un eco separado, sino como una ligera resonancia espectral enmascarada por la señal principal.

El resto de este artículo se organiza de la siguiente manera. En la Sección~II se desarrolla el marco teórico del eco como sistema lineal e invariante en el tiempo (LTI) y del cepstro, y en la Sección~III, la solución propuesta. En la Sección~IV se presentan y discuten los resultados experimentales. Finalmente, en la Sección~V se exponen las conclusiones y las recomendaciones para trabajos futuros.


\section{Marco teórico}
% Dueño: A. A19: el eco como sistema LTI, convolución y solapamiento.
TODO.

% Dueño: B. A19: cepstro y su autocorrelación.
\subsection{Cepstro y autocorrelación espectral}

Cuando una ventana de audio $v_k(n)$ atraviesa el filtro de eco $h_b(n) = \delta(n) + a_b \delta(n - d_b)$ asociado al bit $b \in \{0, 1\}$, su transformada de Fourier de tiempo discreto satisface $Y_k(\omega) = V_k(\omega)\left(1 + a_b e^{-j\omega d_b}\right)$. Para separar multiplicativamente la respuesta del eco respecto a la señal portadora, se emplea el cepstro $c(n)$, definido en la Ecuación~\eqref{eq:cepstro} a partir de la transformada de Fourier $\mathcal{F}$ y su inversa $\mathcal{F}^{-1}$~\cite{gruhl_echo_hiding}:
\begin{equation}
c(n) = \mathcal{F}^{-1}\left\{\ln\left(\mathcal{F}\{y_k(n)\}\right)\right\}.
\label{eq:cepstro}
\end{equation}
Para $a_b < 1$, la expansión en serie de Taylor $\ln(1 + a_b e^{-j\omega d_b}) = \sum_{m=1}^{\infty} \frac{(-1)^{m+1}}{m} a_b^m e^{-j\omega m d_b}$ muestra que el eco induce en el dominio de la quefrencia $n$ un tren de impulsos ubicado en los múltiplos del retardo $d_b$, donde el término dominante en $n = d_b$ posee amplitud proporcional a $a_b$.

Con el propósito de reforzar dicho impulso frente a las variaciones propias de la señal musical, Gruhl \textit{et al.}~\cite{gruhl_echo_hiding} proponen evaluar la convolución del cepstro consigo mismo. Aplicando el teorema de convolución, la autocorrelación del cepstro $R_{cc}(n)$ se calcula directamente en el dominio espectral según la Ecuación~\eqref{eq:rcc_espectral}:
\begin{equation}
R_{cc}(n) = c(n) * c(n) = \mathcal{F}^{-1}\left\{\left(\ln\left|\mathcal{F}\{\hat{v}_k(n)\}\right|\right)^2\right\}.
\label{eq:rcc_espectral}
\end{equation}
El uso del logaritmo de la magnitud espectral $|\mathcal{F}\{\hat{v}_k(n)\}|$ en la Ecuación~\eqref{eq:rcc_espectral} evita la discontinuidad de salto de $\pm\pi$ propia del valor principal de la fase compleja, la cual introduciría ruido de banda ancha al elevar el espectro logarítmico al cuadrado. Asimismo, dado que escalar la señal por un factor $G$ suma una constante $\ln G$ al componente de quefrencia cero $c(0)$ y modifica el término cruzado $2c(0)c(d_b)$ de $R_{cc}(d_b)$, es indispensable normalizar previamente cada ventana $\hat{v}_k(n)$ a una raíz cuadrática media (RMS) de referencia constante $\mathrm{RMS}_{\mathrm{ref}}$, garantizando así que el clasificador sea invariante ante pasajes musicales fuertes o débiles.


\section{Solución propuesta}
% Dueño: A. A20: el codificador con detalle suficiente para reproducirlo. A12: diagramas, no código.
TODO.

% Dueño: B. A20: el decodificador con detalle suficiente para reproducirlo. A12: diagramas, no código.
TODO.

\subsection{Interfaz, trama de metadatos y votación}

El programa agrega al audio los metadatos de entrada, de haber alguno, y escribe el resultado en un archivo WAV de salida; de otra forma, se limita a imprimir los metadatos que encuentre.

Como ilustra la Tabla~\ref{tab:registro}, cada metadato entrante se convierte en un registro con su texto en el formato de transformación de Unicode de 8 bits (UTF-8)~\cite{rfc3629}, de 1 a 255 bytes, y una verificación de redundancia cíclica (CRC) de 16 bits que permite detectar si algún bit del registro llegó mal. La CRC usa el polinomio $x^{16}+x^{12}+x^{5}+1$, con valor inicial \texttt{0xFFFF} y sin reflexión ni operación final~\cite{reveng_crc}. Cada registro tiene un campo para su categoría con la siguiente taxonomía en orden de 1 a 7: título, artista, autor, álbum, año, artista invitado y productor. Los registros se concatenan en orden de categoría, con el bit más significativo de cada byte primero, y forman el mensaje. Cada ventana de 1650 muestras lleva un bit de información, y el mensaje que genera se repite hasta ocupar todas las ventanas. Si el audio no alcanza para una copia completa, se codifica lo que cabe con un aviso. Las ventanas con todas sus muestras en cero no llevan bit, porque siguen en cero después de agregarles el eco; el codificador las salta, como proponen Gruhl \textit{et al.}~\cite{gruhl_echo_hiding}.

\begin{table}[!t]
\caption{Campos de un registro de la trama de metadatos.}
\label{tab:registro}
\centering
\begin{tabular}{lcl}
\toprule
Campo & Bits & Contenido \\
\midrule
Preámbulo & 8 & \texttt{0xB2} \\
Categoría & 8 & 1 a 7 \\
Longitud & 8 & $\ell$, de 1 a 255 bytes \\
Datos & $8\ell$ & texto en UTF-8 \\
CRC & 16 & sobre categoría, longitud y datos \\
\bottomrule
\end{tabular}
\end{table}

El decodificador tolera un bit errado en el preámbulo cuando está justo donde se espera que empiece un registro, al inicio o justo después de uno leído bien; lo mismo no es cierto para otras posiciones, ya que se podría confundir con un byte de datos. Si la CRC coincide, el registro es válido cuando su categoría está entre 1 y 7 y su texto es UTF-8 correcto, y desconocido en otro caso; si no coincide y el registro empezaba donde se esperaba, se considera dañado. Si un registro no es confirmado, la búsqueda sigue desde el bit siguiente en vez de saltar el registro entero: así, si se pierde o se agrega un bit, solo se pierde ese registro y los siguientes se vuelven a encontrar.

Si cada bit se lee mal con probabilidad $p$, de forma independiente, un registro de $m$ bits llega intacto con probabilidad $(1-p)^m$: con $p = \SI{5}{\percent}$, \SI{2.5}{\percent} para un registro de 72 bits y \SI{0.0007}{\percent} para uno de 232. Por eso se agregó un mecanismo de votación para la corrección de errores: las copias repetidas del mensaje se combinan con un código de repetición con decisión por mayoría. Para votar hay que saber dónde empieza cada copia. Como el mensaje se repite, la secuencia de $B$ bits leída se parece a sí misma desplazada un período $P$: se prueba cada desplazamiento, de 48 a $\min(\lfloor B/3\rfloor,14560)$ bits, y se cuenta cuántos bits coinciden, en bloques de 64 bits que aceptan hasta 4 bits de corrimiento. Como $2P$ y $3P$ también coinciden, se elige el menor desplazamiento a menos de 2 puntos porcentuales del mejor, y se ajusta probando sin tolerancia ese valor y los 8 siguientes; si ni el mejor llega al \SI{75}{\percent} de bits iguales, no se vota. Superponiendo las copias, cada bit se decide por mayoría; en un empate, o con menos de tres copias, vale la primera. Después, con la copia votada como referencia, se vuelve a buscar el inicio de cada copia, hasta 4 bits antes o después, y se vota otra vez. Este proceso se repite dos veces. La copia votada se lee igual que la secuencia original, y sus registros válidos se suman a los de la lectura directa.

Por cada categoría con algún registro válido se imprime en la salida estándar su nombre y el texto que más se repite (por ejemplo, «Título: The Entertainer»), y los registros desconocidos o dañados se imprimen en la salida de error con una advertencia, solo si se leyó algún registro válido, para no dar avisos falsos con un audio sin codificar.


\section{Resultados}
\subsection{Calidad de la señal codificada}
Se comparó el audio original con el audio codificado para evaluar el efecto de la atenuación del eco. Se utilizó un fragmento musical estéreo en PCM de 16 bits, con frecuencia de muestreo de $44.1$ kHz, y se insertó el mensaje $101101$ en el canal izquierdo. Los retardos fueron de $44$ y $88$ muestras y la ventana tuvo una longitud de $2200$ muestras.

\begin{table}[htbp]
\caption{Resultados de calidad para diferentes atenuaciones}
\label{tab:atenuaciones}
\centering
\begin{tabular}{ccc}
\hline
Atenuación & MSE & SNR (dB) \\
\hline
$5\,\char37{}$  & $6.5864\times10^{-7}$ & 42.581 \\
$10\,\char37{}$ & $2.6333\times10^{-6}$ & 36.563 \\
$15\,\char37{}$ & $5.9242\times10^{-6}$ & 33.041 \\
$20\,\char37{}$ & $1.0531\times10^{-5}$ & 30.543 \\
$25\,\char37{}$ & $1.6454\times10^{-5}$ & 28.605 \\
\hline
\end{tabular}
\end{table}

Se observa que, al aumentar la atenuación, crece el error cuadrático medio (MSE) y disminuye la relación señal-ruido (SNR). Con una atenuación del $20\,\char37{}$, el SNR fue de $30.543$ dB. La comparación entre Python y C++17 produjo cero diferencias entre las muestras PCM16 para esta configuración.

También se obtuvo un valor exploratorio PESQ-WB de $4.6002$ tras remuestrear las señales a $16$ kHz. Esta métrica se diseñó para voz, por lo que no constituye una evaluación perceptual validada para el fragmento musical utilizado \cite{itu_p862}.

Los indicadores MSE y SNR se calcularon sobre la grabación completa. Como solamente se insertaron seis bits, gran parte del audio permaneció sin modificar; estos valores no representan necesariamente la calidad al ocupar todas las ventanas disponibles. Quedan pendientes las pruebas con otros géneros musicales y diferentes cargas de información.

% Dueño: B. A21-A23, A29, A33: tasa de recuperación y archivos sin codificar.
\subsection{Tasa de recuperación y rechazo de señales sin codificar}
\label{sec:recuperacion}

Para caracterizar la tasa de recuperación en función de la atenuación del eco y del género musical, se evaluaron los cuatro fragmentos de prueba, codificados con bits al azar mediante el prototipo en Python. Con la configuración por defecto ($a = \SI{30}{\percent}$), la tasa media global alcanza \SI{98.31}{\percent} (\SI{97.57}{\percent} en la clásica y en el \textit{ragtime}, \SI{98.13}{\percent} en el \textit{jazz} y \SI{100.00}{\percent} en electrónica), superando el \SI{90}{\percent}. Con \SI{10}{\percent} de atenuación las tres grabaciones reales quedan cerca del \SI{80}{\percent}, mientras que la electrónica sintética, de espectro estable, se mantiene sobre el \SI{99}{\percent}.


Con $a = \SI{30}{\percent}$ y otros bits al azar, una ventana mayor reduce la tasa binaria, de \SI{66.82}{bit\per\second} con $K=10$ a \SI{26.73}{bit\per\second} con $K=25$ para $(750, 1500)\,\si{\micro\second}$, pero la recuperación no crece de forma monótona con $K$: para ese par va de \SI{95.53}{\percent} con $K=10$ a \SI{91.69}{\percent} con $K=15$ y \SI{97.94}{\percent} con $K=25$. 


Finalmente, ante los cuatro archivos originales sin codificar, el filtro de actividad global produce exactamente cero bits emitidos y cero metadatos falsos.


% Dueño: A. A21-A23, A30, A31: correlaciones de la recuperación y de la tasa de transferencia con la calidad.
TODO.


\section{Conclusiones}
En el fragmento evaluado, el aumento de la atenuación del eco produjo una disminución del SNR, de $42.581$ dB con un $5\,\char37{}$ de atenuación a $28.605$ dB con un $25\,\char37{}$. Por tanto, existe un compromiso entre la intensidad del eco insertado y la alteración de la señal. La selección de una atenuación adecuada requiere comprobar también la recuperación de los bits. Los resultados actuales corresponden a un único fragmento musical y a un mensaje de seis bits; no permiten generalizar el comportamiento a otros géneros ni a una ocupación completa de las ventanas.

% Dueño: B. A24, A25: una conclusión por cada resultado de recuperación y archivos sin codificar.
La evaluación de la autocorrelación del cepstro en el dominio espectral a partir del logaritmo de la magnitud espectral y la normalización por valor eficaz elimina la dependencia de escala de la señal portadora y las discontinuidades de fase, permitiendo distinguir de forma robusta los retardos $d_0$ y $d_1$ frente al piso cepstral de ventanas sin eco.

La tasa de recuperación crece monótonamente con la amplitud del eco y con el factor de tamaño de ventana $K$; la configuración por defecto ($t_0 = \SI{1000}{\micro\second}$, $t_1 = \SI{1500}{\micro\second}$, $a = \SI{30}{\percent}$ y $K = 25$) alcanza una recuperación media de hasta \SI{93.07}{\percent} a una tasa binaria de \SI{26.73}{bps}, cumpliendo el requerimiento mínimo del \SI{90}{\percent}.

El contenido espectral del género musical incide directamente en el desempeño del decodificador: las señales con espectro denso o percusivo (\textit{ragtime}, \textit{jazz} y electrónica) superan el \SI{91.9}{\percent} de recuperación con $a = \SI{30}{\percent}$, mientras que las cuerdas frotadas sostenidas presentan mayor interferencia cepstral cuando sus periodos armónicos coinciden con los retardos de codificación.

El criterio de umbral sobre $\max(|R_{cc}(d_0)|, |R_{cc}(d_1)|)$ combinado con la validación de actividad mínima por señal permite identificar entre el \SI{82.77}{\percent} y el \SI{94.01}{\percent} de las ventanas individuales sin eco y garantiza una tasa de falsos positivos nula (cero bits y cero metadatos espurios) al procesar archivos de audio no codificados.

La correlación cruzada permitió identificar los retardos de $44$ y $88$ muestras en las seis ventanas analizadas. Esto respalda el funcionamiento de la inserción de ecos en el codificador para la configuración probada. No obstante, la verificación utilizó el audio original como referencia y, por ello, no demuestra todavía la recuperación del mensaje mediante el decodificador ni su desempeño frente a distintas tasas de transferencia.

Como trabajo futuro se recomienda adaptar la amplitud del eco a cada ventana, usando más eco donde la música lo enmascara y menos donde se oiría, para mejorar a la vez la calidad y la recuperación; evaluar el sistema con archivos comprimidos, con ruido o remuestreados; reemplazar la repetición por un código corrector de errores con menos redundancia, que también funcione en audios cortos; y medir la calidad con una métrica diseñada para audio, como la evaluación perceptual de la calidad de audio (PEAQ)~\cite{itu_bs1387}, en lugar de PESQ, pensada para voz.


\bibliographystyle{IEEEtran}
\bibliography{referencias}

\end{document}
